% !TEX root = ../main.tex
\begin{ouverture}[titrage]{3}{Titrages pH-métriques et conductimétriques}{Partie 1 · Constitution et transformations de la matière}
{Titre massique\enspace·\enspace Équivalence\enspace·\enspace Suivi pH-métrique\enspace·\enspace Suivi conductimétrique}
\objectif{Préparer une solution de concentration donnée à partir d'une solution commerciale.}
\objectif{Établir la composition du système au cours d'un titrage, avant et après l'équivalence.}
\objectif{Exploiter un titrage pour déterminer une quantité de matière, une concentration ou une masse.}
\objectif{Justifier l'évolution de la pente d'une courbe de titrage conductimétrique.}
\objectif{Représenter l'évolution des quantités de matière avec un programme Python.}
\prerequis{Réactions acide-base et couples (chapitre 1).}
\prerequis{Conductivité et loi de Kohlrausch (chapitre 2).}
\prerequis{Titrage avec suivi colorimétrique, équivalence, tableau d'avancement.}
\end{ouverture}
\essentiel
\essSection{1}{Préparer une solution à partir d'une solution commerciale}
\begin{definition}[Titre massique et densité]
Le titre massique d'une solution est $t=\dfrac{m_\text{soluté}}{m_\text{solution}}$ (souvent exprimé en \%).
Sa densité est $d=\dfrac{\rho}{\rho_\text{eau}}$, avec $\rho_\text{eau}=\qty{1,00}{kg.L^{-1}}$. La concentration en
soluté apporté s'en déduit :
\[c=\frac{t\times\rho}{M}=\frac{t\times d\times\rho_\text{eau}}{M}.\]
\end{definition}
\begin{exemple}
Une solution commerciale d'acide chlorhydrique à 37 \% a une densité de 1,18. Un litre a une masse de 1\,180 g et
contient $0{,}37\times1\,180\approx\qty{437}{g}$ de chlorure d'hydrogène, soit
$c=\frac{437}{36{,}5}\approx\qty{12}{mol.L^{-1}}$.
\end{exemple}
\begin{attention}
Les solutions commerciales sont souvent corrosives : lunettes, gants et blouse, manipulation sous hotte.
Pour diluer un acide concentré, on verse l'acide dans l'eau, et jamais l'inverse.
\end{attention}
\essSection{2}{Principe d'un titrage}
\begin{definition}[Titrage et équivalence]
Titrer une espèce, c'est déterminer sa quantité de matière grâce à une réaction support \textbf{rapide},
\textbf{totale} et \textbf{unique} avec une espèce titrante de concentration connue. L'\textbf{équivalence} est
atteinte lorsque les réactifs ont été introduits dans les proportions stœchiométriques : le réactif limitant
change. Pour la réaction $a\,\mathrm{A}+b\,\mathrm{B}\longrightarrow\dots$, avec A titré et B titrant :
\[\frac{c_\mathrm{A}\,V_\mathrm{A}}{a}=\frac{c_\mathrm{B}\,V_\mathrm{E}}{b},\]
où $V_\mathrm{E}$ est le volume de solution titrante versé à l'équivalence.
\end{definition}
\schema{03-montage-titrage}
\begin{attention}
Le volume à l'équivalence $V_\mathrm{E}$ n'a rien à voir avec un état d'équilibre : à l'équivalence, la réaction
support, totale, a consommé entièrement les deux réactifs. L'eau distillée ajoutée dans le bécher, par exemple
pour immerger la sonde, ne modifie pas la quantité d'espèce titrée, donc pas $V_\mathrm{E}$.
\end{attention}
\essSection{3}{Composition du système au cours du titrage}
\begin{propriete}[Avant et après l'équivalence]
Avant l'équivalence, l'espèce titrante versée est entièrement consommée : l'espèce titrée est en excès.
Après l'équivalence, l'espèce titrée a disparu et l'espèce titrante versée s'accumule. Les ions spectateurs
(\ce{Na+}, \ce{Cl-}…) ne réagissent pas : leur quantité ne dépend que de ce qui a été introduit.
\end{propriete}
\begin{codecote}[0.58]
\begin{pythoncol}
import matplotlib.pyplot as plt
n0, cB = 0.20, 0.100  # n0 : mmol, cB : mol/L
V = [k * 0.1 for k in range(41)]   # en mL
nH, nHO, nNa = [], [], []
for v in V:
    nb = cB * v       # ions HO- versés (mmol)
    nH.append(max(n0 - nb, 0))
    nHO.append(max(nb - n0, 0))
    nNa.append(nb)
plt.plot(V, nH, label="H3O+")
plt.plot(V, nHO, label="HO-")
plt.plot(V, nNa, label="Na+")
plt.legend()
plt.show()
\end{pythoncol}
\cote
\fontsize{9.5}{12.5}\selectfont
\capacite{Capacité numérique}\par\smallskip
Titrage de \qty{0,20}{mmol} d'ions oxonium par la soude. La réaction support \ce{H3O+ + HO- -> 2 H2O} étant
totale, \code{max(n0 - nb, 0)} donne la quantité d'ions oxonium restante et \code{max(nb - n0, 0)} celle
des ions hydroxyde en excès. Le volume équivalent correspond au point où les deux courbes s'annulent ensemble.
\end{codecote}
\schema{03-python-quantites}
\essSection{4}{Titrage avec suivi pH-métrique}
\begin{propriete}[Courbe de titrage pH-métrique]
Au voisinage de l'équivalence, le pH varie brutalement : c'est le saut de pH. On détermine $V_\mathrm{E}$ par la
méthode des \textbf{tangentes parallèles} (le point d'équivalence E est sur la droite équidistante des deux
tangentes) ou par la méthode de la \textbf{dérivée} ($V_\mathrm{E}$ correspond au maximum de
$\left|\mathrm{dpH}/\mathrm{d}V\right|$). Le pH à l'équivalence vaut 7 pour un acide fort titré par une base forte ;
il est supérieur à 7 pour un acide faible, comme l'acide éthanoïque, titré par une base forte.
\end{propriete}
\schema{03-courbes-ph}
\begin{savoirfaire}[Choisir un indicateur coloré]
Pour un titrage colorimétrique, on choisit un indicateur dont la zone de virage contient le pH à l'équivalence :
hélianthine (3,1 -- 4,4), bleu de bromothymol (6,0 -- 7,6), phénolphtaléine (8,2 -- 10,0).
Ainsi, la phénolphtaléine convient pour titrer l'acide éthanoïque par la soude ; ce point est approfondi au chapitre 9.
\end{savoirfaire}
\essSection{5}{Titrage avec suivi conductimétrique}
\begin{propriete}[Courbe de titrage conductimétrique]
Si l'on ajoute assez d'eau pour que la dilution soit négligeable, la conductivité varie de façon affine avec le
volume versé, avec une pente différente avant et après l'équivalence : $V_\mathrm{E}$ est l'abscisse du point
d'intersection des deux droites. Le signe et la valeur de chaque pente s'expliquent par les ions qui apparaissent
et disparaissent, et par leurs conductivités molaires ioniques. Ce suivi convient aussi à des réactions qui ne sont
pas acide-base, comme la précipitation \ce{Ag+ + Cl- -> AgCl(s)}.
\end{propriete}
\schema{03-courbe-sigma}
\begin{savoirfaire}[Justifier l'évolution de la pente]
Écrire l'équation de la réaction support, puis faire le bilan des ions : pour chaque ion titrant versé, quels ions
sont ajoutés, lesquels sont retirés ? Comparer leurs conductivités molaires : le remplacement d'un ion par un ion
moins conducteur fait diminuer $\sigma$ ; l'accumulation d'ions fait croître $\sigma$.
\end{savoirfaire}
\begin{rediger}
Écrire « à l'équivalence, les réactifs ont été introduits dans les proportions stœchiométriques de la réaction
support » et non « les réactifs sont en quantités égales ». Citer la relation à l'équivalence avec ses notations,
puis la méthode utilisée pour déterminer $V_\mathrm{E}$ ; terminer par une phrase qui répond à la question posée.
\end{rediger}
\exercices

\rubrique{Découverte}
\begin{blocLivre}{ex-3-1}
\begin{exercice}[1]{Une solution commerciale}
Une solution commerciale d'acide chlorhydrique porte les indications : « 23 \% en masse ; densité 1,11 ».
\begin{donnees}
$M(\ce{HCl})=\qty{36,5}{g.mol^{-1}}$ ; $\rho_\text{eau}=\qty{1,00}{kg.L^{-1}}$.
\end{donnees}
\begin{questions}
\item Calculer la masse d'un litre de solution, puis la masse de chlorure d'hydrogène qu'il contient.
\item En déduire la concentration en acide apporté de la solution commerciale.
\item Quel volume de solution commerciale faut-il prélever pour préparer 500,0 mL de solution à \qty{0,100}{mol.L^{-1}} ?
\item Citer deux précautions de sécurité pour cette préparation.
\end{questions}
\end{exercice}
\end{blocLivre}
\begin{blocLivre}{ex-3-2}
\begin{exercice}[1]{Exploiter l'équivalence}
Pour chaque titrage, écrire l'équation de la réaction support, la relation à l'équivalence, puis calculer la
concentration de l'espèce titrée.
\begin{questions}
\item 20,0 mL d'acide chlorhydrique sont titrés par la soude à \qty{0,100}{mol.L^{-1}} ; $V_\mathrm{E}=\qty{12,6}{mL}$.
\item 10,0 mL d'une solution d'ammoniac \ce{NH3} sont titrés par l'acide chlorhydrique à \qty{0,0500}{mol.L^{-1}} ;
$V_\mathrm{E}=\qty{14,2}{mL}$.
\end{questions}
\end{exercice}
\end{blocLivre}
\begin{blocLivre}{ex-3-3}
\begin{exercice}[1]{Lire une courbe de titrage pH-métrique}
On titre $V_\mathrm{A}=\qty{20,0}{mL}$ d'acide chlorhydrique par la soude à $c_\mathrm{B}=\qty{1,00e-2}{mol.L^{-1}}$.
\schema{03-ex03}
\begin{questions}
\item Déterminer le volume à l'équivalence par la méthode des tangentes parallèles. Relever le pH à l'équivalence.
\item En déduire la concentration $c_\mathrm{A}$ de l'acide chlorhydrique.
\item Le pH initial est-il cohérent avec cette concentration ? Justifier.
\end{questions}
\end{exercice}
\end{blocLivre}
\begin{blocLivre}{ex-3-4}
\begin{exercice}[1]{Lire une courbe de titrage conductimétrique}
On titre 20,0 mL de soude, additionnés de 180 mL d'eau, par l'acide chlorhydrique à \qty{0,0500}{mol.L^{-1}}.
\schema{03-ex04}
\begin{donnees}
$\lambda$ (mS$\cdot$m$^2\cdot$mol$^{-1}$) : \ce{HO-} : 19,9 ; \ce{Cl-} : 7,63 ; \ce{H3O+} : 35,0 ; \ce{Na+} : 5,01.
\end{donnees}
\begin{questions}
\item Déterminer le volume à l'équivalence, puis la concentration de la soude.
\item Expliquer pourquoi la conductivité diminue avant l'équivalence, puis augmente plus vite après.
\end{questions}
\end{exercice}
\end{blocLivre}
\begin{blocLivre}{ex-3-5}
\begin{exercice}[1]{Choisir la verrerie d'un titrage}
On dispose de : bécher de 100 mL ; éprouvette graduée de 25 mL ; pipette jaugée de 20,0 mL ; burette graduée de
25 mL ; fiole jaugée de 100,0 mL ; pissette d'eau distillée.
\begin{questions}
\item Quelle verrerie choisir pour prélever la solution titrée ? Pour verser la solution titrante ? Justifier.
\item Pourquoi rince-t-on la burette avec la solution titrante avant de la remplir ?
\item On ajoute de l'eau distillée dans le bécher pour immerger la sonde. Cela modifie-t-il le volume à l'équivalence ?
\end{questions}
\end{exercice}
\end{blocLivre}
\rubrique{Entraînement}
\begin{blocLivre}{ex-3-6}
\begin{exercice}[2]{La composition du système au cours d'un titrage}
On titre 20,0 mL d'acide chlorhydrique à \qty{1,0e-2}{mol.L^{-1}} par la soude à \qty{1,0e-2}{mol.L^{-1}}.
\begin{questions}
\item Écrire l'équation de la réaction support et calculer le volume à l'équivalence.
\item Dresser un tableau donnant les quantités de matière des ions \ce{H3O+}, \ce{Cl-}, \ce{Na+} et \ce{HO-}
pour $V_\mathrm{B}=\qty{10,0}{mL}$, $\qty{20,0}{mL}$ et $\qty{30,0}{mL}$.
\item Calculer le pH du mélange pour $V_\mathrm{B}=\qty{10,0}{mL}$ et pour $V_\mathrm{B}=\qty{30,0}{mL}$
(on admet que $[\ce{H3O+}]\,[\ce{HO-}]=10^{-14}$ en mol$^2\cdot$L$^{-2}$ à 25 °C).
\end{questions}
\end{exercice}
\end{blocLivre}
\begin{blocLivre}{ex-3-7}
\begin{exercice}[2]{Interpréter les pentes d'une courbe conductimétrique}
On compare deux titrages par la soude : celui d'une solution d'acide chlorhydrique et celui d'une solution d'acide
éthanoïque, de mêmes concentrations, additionnées d'un grand volume d'eau.
\begin{donnees}
$\lambda$ (mS$\cdot$m$^2\cdot$mol$^{-1}$) : \ce{H3O+} : 35,0 ; \ce{HO-} : 19,9 ; \ce{Na+} : 5,01 ; \ce{CH3COO-} : 4,09 ;
l'acide éthanoïque est très peu dissocié dans l'eau.
\end{donnees}
\begin{questions}
\item Pour chaque titrage, écrire l'équation de la réaction support.
\item Pour l'acide chlorhydrique, faire le bilan des ions ajoutés et retirés pour une mole d'ions \ce{HO-} versée,
avant puis après l'équivalence. En déduire le signe des pentes.
\item Même question pour l'acide éthanoïque. Pourquoi la conductivité augmente-t-elle dès le début du titrage ?
\item Prévoir le rapport des pentes après et avant l'équivalence dans le cas de l'acide éthanoïque. Esquisser la courbe.
\end{questions}
\end{exercice}
\end{blocLivre}
\begin{blocLivre}{ex-3-8}
\begin{exercice}[2]{Titrer un vinaigre}
Un vinaigre est annoncé à 6,0 \% d'acidité (6,0 g d'acide éthanoïque pour 100 g de vinaigre). On le dilue dix fois,
puis on titre 10,0 mL de solution diluée par la soude à \qty{0,100}{mol.L^{-1}}.
\schema{03-ex08}
\begin{donnees}
$M(\ce{CH3COOH})=\qty{60,05}{g.mol^{-1}}$ ; masse volumique du vinaigre : \qty{1,01}{g.mL^{-1}}.
\end{donnees}
\begin{questions}
\item Écrire l'équation de la réaction support du titrage.
\item Déterminer le volume à l'équivalence et le pH à l'équivalence. Pourquoi ce pH est-il supérieur à 7 ?
\item En déduire la concentration en acide éthanoïque du vinaigre, puis son pourcentage d'acidité. Conclure.
\item Quel indicateur coloré pourrait-on utiliser pour un titrage colorimétrique ? Justifier.
\end{questions}
\end{exercice}
\end{blocLivre}
\begin{blocLivre}{ex-3-9}
\begin{exercice}[2]{Les ions chlorure d'un sérum physiologique}
On dilue un sérum physiologique dix fois. On titre 20,0 mL de solution diluée, additionnés de 200 mL d'eau, par une
solution de nitrate d'argent (\ce{Ag+}, \ce{NO3-}) à \qty{0,0500}{mol.L^{-1}}. Il se forme un précipité blanc de
chlorure d'argent.
\schema{03-ex09}
\begin{donnees}
$\lambda$ (mS$\cdot$m$^2\cdot$mol$^{-1}$) : \ce{Cl-} : 7,63 ; \ce{NO3-} : 7,14 ; \ce{Ag+} : 6,19 ; \ce{Na+} : 5,01 ;
$M(\ce{NaCl})=\qty{58,44}{g.mol^{-1}}$.
\end{donnees}
\begin{questions}
\item Écrire l'équation de la réaction support. Pourquoi ce titrage ne peut-il pas être suivi par pH-métrie ?
\item Justifier la faible diminution de la conductivité avant l'équivalence, puis son augmentation.
\item Déterminer le volume à l'équivalence, puis la concentration en ions chlorure du sérum.
\item Le sérum est annoncé à 9,0 g$\cdot$L$^{-1}$ de chlorure de sodium. Conclure.
\end{questions}
\end{exercice}
\end{blocLivre}
\begin{blocLivre}{ex-3-10}
\begin{exercice}[2]{Programmer l'évolution des quantités de matière}
On titre 10,0 mL d'acide éthanoïque à \qty{0,100}{mol.L^{-1}} par la soude à \qty{0,100}{mol.L^{-1}} ; la réaction
support est considérée comme totale.
\begin{python}
n0, cB = ..., 0.100        # n0 en mmol
V = [k * 0.5 for k in range(41)]
nA, nAm, nHO = [], [], []
for v in V:
    nb = cB * v
    nA.append(...)           # acide éthanoïque restant
    nAm.append(min(nb, n0))  # ion éthanoate formé
    nHO.append(...)          # ions hydroxyde en excès
\end{python}
\begin{questions}
\item \capacite{Capacité numérique}Écrire l'équation de la réaction support et calculer $n_0$, en mmol.
\item Compléter les deux lignes incomplètes du programme.
\item Pourquoi la quantité d'ions éthanoate cesse-t-elle d'augmenter ? Pour quel volume ?
\item Représenter l'allure des courbes obtenues, ainsi que celle de la quantité d'ions \ce{Na+}.
\end{questions}
\end{exercice}
\end{blocLivre}
\begin{blocLivre}{ex-3-11}
\begin{exercice}[2]{L'acidité d'un lait}
L'acidité d'un lait s'exprime en degrés Dornic (°D) : 1 °D correspond à 0,1 g d'acide lactique par litre de lait.
Un lait est considéré comme frais si son acidité ne dépasse pas 18 °D. On titre 10,0 mL de lait par une solution
d'hydroxyde de sodium dite « soude Dornic », de concentration $c=\frac19$ mol$\cdot$L$^{-1}$, en présence de
phénolphtaléine. Le virage est obtenu pour $V_\mathrm{E}=\qty{1,6}{mL}$.
\begin{donnees}
Acide lactique \ce{C3H6O3}, noté AH : $M=\qty{90,08}{g.mol^{-1}}$ ; \ce{AH + HO- -> A- + H2O}.
\end{donnees}
\begin{questions}
\item Calculer la quantité, puis la masse d'acide lactique contenue dans 10,0 mL de lait.
\item En déduire l'acidité du lait en degrés Dornic. Le lait est-il frais ?
\item Montrer que, pour une prise d'essai de 10,0 mL, l'acidité en °D vaut dix fois le volume équivalent exprimé en mL.
\end{questions}
\end{exercice}
\end{blocLivre}
\rubrique{Objectif bac}
\begin{blocLivre}{ex-3-12}
\begin{exercice}[3]{Contrôle d'un déboucheur}
\begin{docu}{1}{Étiquette d'un déboucheur liquide}
Hydroxyde de sodium : 10 \% en masse ; densité : 1,10. Danger : provoque de graves brûlures de la peau et des lésions
oculaires.
\end{docu}
\begin{docu}{2}{Protocole et résultat}
On dilue le déboucheur 25 fois. On titre 20,0 mL de solution diluée par l'acide chlorhydrique à
\qty{0,100}{mol.L^{-1}}, avec suivi pH-métrique.
\end{docu}
\schema{03-ex12}
\begin{donnees}
$M(\ce{NaOH})=\qty{40,0}{g.mol^{-1}}$ ; matériel : pipettes jaugées de 5,0, 10,0 et 20,0 mL ; fioles jaugées de 100,0
et 250,0 mL.
\end{donnees}
\begin{questions}
\item Proposer un protocole de dilution et citer les précautions de sécurité nécessaires.
\item Écrire l'équation de la réaction support du titrage et déterminer le volume à l'équivalence.
\item \initiative L'indication de l'étiquette est-elle vérifiée ?
\item Quel volume à l'équivalence aurait-on obtenu si le titre massique valait exactement 10 \% ?
\end{questions}
\end{exercice}
\end{blocLivre}
\begin{blocLivre}{ex-3-13}
\begin{exercice}[3]{Titrage d'un comprimé d'aspirine}
\begin{docu}{1}{Protocole}
On broie un comprimé d'aspirine « 500 » et on dissout la poudre dans un mélange d'eau et d'éthanol, complété à
200,0 mL. On titre 50,0 mL de cette solution par la soude à \qty{0,0500}{mol.L^{-1}}, à froid. Un logiciel trace
le pH et sa dérivée en fonction du volume versé.
\end{docu}
\schema{03-ex13}
\begin{donnees}
Acide acétylsalicylique \ce{C9H8O4}, noté AH (une seule fonction acide titrée) : $M=\qty{180,16}{g.mol^{-1}}$.
À chaud, l'ion hydroxyde réagit aussi avec la fonction ester de la molécule.
\end{donnees}
\begin{questions}
\item Écrire l'équation de la réaction support du titrage.
\item Déterminer le volume à l'équivalence. Quelle méthode est la plus précise ici ? Justifier.
\item \initiative Le comprimé contient-il bien 500 mg d'acide acétylsalicylique ?
\item Pourquoi le titrage doit-il être réalisé à froid ? Dans quel sens un titrage à chaud fausserait-il le résultat ?
\end{questions}
\end{exercice}
\end{blocLivre}
\begin{blocLivre}{ex-3-14}
\begin{exercice}[3]{Un vinaigre frelaté ?}
\begin{docu}{1}{Une fraude possible}
Ajouter de l'acide chlorhydrique à un vinaigre trop faible augmente son acidité à moindre coût ; cette pratique est
interdite. Pour contrôler un vinaigre suspect, on le dilue vingt fois : on obtient la solution S. On titre 20,0 mL de
S, additionnés de 180 mL d'eau, par la soude à \qty{0,100}{mol.L^{-1}}, avec suivi conductimétrique.
\end{docu}
\schema{03-ex14}
\begin{donnees}
$\lambda$ (mS$\cdot$m$^2\cdot$mol$^{-1}$) : \ce{H3O+} : 35,0 ; \ce{HO-} : 19,9 ; \ce{Na+} : 5,01 ; \ce{CH3COO-} : 4,09 ;
\ce{Cl-} : 7,63. Dans un mélange, l'acide fort réagit en premier avec les ions hydroxyde.
\end{donnees}
\begin{questions}
\item Écrire les équations des deux réactions supports, dans l'ordre où elles se produisent.
\item Interpréter l'allure de chacune des trois portions de la courbe.
\item \initiative Déterminer les concentrations des deux acides dans le vinaigre. Le vinaigre est-il frelaté ?
\end{questions}
\end{exercice}
\end{blocLivre}
\begin{blocLivre}{ex-3-15}
\begin{exercice}[3]{Quand la dilution n'est plus négligeable}
On titre 20,0 mL d'acide chlorhydrique de concentration inconnue par la soude à \qty{0,100}{mol.L^{-1}}, sans ajouter
d'eau. Les points ne s'alignent plus sur deux droites.
\schema{03-ex15}
\begin{questions}
\item À l'aide de la loi de Kohlrausch, montrer que $\sigma\,(V_0+V_\mathrm{B})$ est une fonction affine de $V_\mathrm{B}$
de part et d'autre de l'équivalence, avec $V_0=\qty{20,0}{mL}$.
\item \capacite{Capacité numérique}On calcule la conductivité corrigée $\sigma'=\sigma\,(V_0+V_\mathrm{B})/V_0$.
Quelle valeur prend-elle pour $V_\mathrm{B}=\qty{10,0}{mL}$, où l'on mesure $\sigma=\qty{4,21}{mS.cm^{-1}}$ ?
\item Avec la conductivité corrigée, on obtient deux droites qui se coupent pour $V_\mathrm{B}=\qty{10,0}{mL}$.
En déduire la concentration de l'acide chlorhydrique.
\item \initiative Un élève trace des droites passant au plus près des cinq premiers et des cinq derniers points mesurés,
sans correction. Elles se coupent pour \qty{7,9}{mL}. Quelle erreur relative commet-il ? Comment éviter le problème
sans calcul ?
\end{questions}
\end{exercice}
\end{blocLivre}
